\documentclass[12pt]{article}
\usepackage[utf8]{inputenc}
\usepackage[T1]{fontenc}
\usepackage{amsmath,amssymb,amsthm}
\usepackage{geometry}
\geometry{margin=1in}
\usepackage{enumitem}
\setlist{nosep}

\linespread{1.1} % Slightly increased line spacing for better readability
% -------------------- Macros --------------------
\newcommand{\N}{\mathbb{N}}
\newcommand{\Z}{\mathbb{Z}}
\newcommand{\Q}{\mathbb{Q}}
\newcommand{\R}{\mathbb{R}}
\newcommand{\C}{\mathbb{C}}
\newcommand{\K}{\mathbb{K}} % used to denote either R or C
\newcommand{\eps}{\varepsilon}
\newcommand{\dd}{\,\mathrm{d}}

\newcommand{\NumberedDefinition}[3]{% #1 number, #2 title, #3 body
\paragraph*{Definition #1 ( #2 ).} #3\par}
\newcommand{\NumberedTheorem}[3]{% #1 number, #2 title, #3 body
\paragraph*{Theorem #1 ( #2 ).} #3\par}
\newcommand{\NumberedProposition}[3]{% #1 number, #2 title, #3 body
\paragraph*{Proposition #1 ( #2 ).} #3\par}
\newcommand{\NumberedLemma}[3]{% #1 number, #2 title, #3 body
\paragraph*{Lemma #1 ( #2 ).} #3\par}
\newcommand{\NumberedCorollary}[3]{% #1 number, #2 title, #3 body
\paragraph*{Corollary #1 ( #2 ).} #3\par}
\newcommand{\NumberedRemark}[3]{% #1 number, #2 title, #3 body
\paragraph*{Remark #1 ( #2 ).} #3\par}
\newcommand{\NumberedExample}[3]{% #1 number, #2 title, #3 body
\paragraph*{Example #1 ( #2 ).} #3\par}

% This chapter translation: Original German "Die Exponentialfunktion" (The Exponential Function)
% Numbering (Definition/Satz) preserved exactly as in source. Do NOT renumber.

\begin{document}
\section*{Chapter 5 \\ Differential Calculus of One Variable}

\subsection*{5.1 Differentiability}

In the following, unless otherwise stated, let the underlying field $\K$ be either $\R$ or $\C$ and $(Y, \|\cdot\|_Y)$ a normed $\K$-vector space, $M \subset \K$, $M \neq \emptyset$, $x_0 \in M$ be an accumulation point (HP)\footnotemark[1] of $M$ and $f : M \to Y$ a function.

\footnotetext[1]{For $M \subset Y$ a point $y \in Y$ is called \textit{accumulation point} of $M$, if for every neighborhood $U$ of $y$ in $Y$ holds $(M \cap U) \setminus \{y\} \neq \emptyset$.}

\NumberedDefinition{5.1.1}{}{The function $f : M \to Y$ is called \textit{differentiable} at the point $x_0 \in M$, if the limit\footnotemark[2]
\[
\lim_{x \to x_0} \frac{f(x) - f(x_0)}{x - x_0}
\]
exists in $Y$. If $f$ is differentiable in $x_0$, we use the notations
\[
f'(x_0) := Df(x_0) := \frac{df}{dx}(x_0) := \lim_{x \to x_0} \frac{f(x) - f(x_0)}{x - x_0} := \lim_{h \to 0} \frac{f(x_0 + h) - f(x_0)}{h}
\]
and call all these quantities \textit{derivative} of $f$ in $x_0$.}

\footnotetext[2]{Let $x_0 \in \K$ and assume there exists a sequence $(x_n)_n \subset M$ with $\lim_{n \to \infty} x_n = x_0$. Then we define:
\[
y := \lim_{x \to x_0} f(x),
\]
if for every sequence $(x_n)_n \subset M$ with $x_n \to x_0$ for $n \to \infty$ holds:
\[
f(x_n) \xrightarrow{n\to\infty} y.
\]}

\NumberedExample{5.1.2}{}{
(1) $f : \C \to \C$, $f(z) = z^n$, $n \in \N_{\ge 1}$ is differentiable for all $z \in \C$ and it holds:
\[
f'(z) = nz^{n-1}
\]
(2) $f : \C \to \C, f(z) := e^z$ is differentiable for all $z \in \C$ and $f'(z) = e^z$.
\par
(3) $f : \R \setminus \{0\} \to \R, f(x) = \frac{1}{x}$ is differentiable for every $x \in \R \setminus \{0\}$ and $f'(x) = -\frac{1}{x^2}$.
\par
(4) $f : \R \to \R^2, f(x) = (x, x^2)^\mathsf{T}$ is differentiable for $x \in \R$ and $f'(x) = (1, 2x)^T$.
\par
(5) $f : \R \to \R_{\ge 0}, f(x) = |x|$ is not differentiable in $x_0 = 0$.
\par
(6) $f : \C \to \R, f(z) = \operatorname{Re}(z)$ is not differentiable in any point $z \in \C$.
}

\paragraph*{Proof.}
(1)
\[
\frac{z^n - z_0^n}{z - z_0} = \sum_{j=0}^{n-1} z^j z_0^{n-1-j} \xrightarrow{z \to z_0} \sum_{j=0}^{n-1} z_0^j z_0^{n-1-j} = z_0^{n-1} \sum_{j=0}^{n-1} 1 = n z_0^{n-1}
\]

(2)
\[
\frac{e^z - e^{z_0}}{z - z_0} = e^{z_0} \frac{e^{z-z_0} - 1}{z - z_0}.
\]
In [8, Ex. 4.3.14(2)] it was shown that the fraction on the right-hand side converges to 1 for $z \to z_0$, so that the assertion follows.

(3)
\[
\frac{\frac{1}{x_0+h} - \frac{1}{x_0}}{h} = \frac{-h}{x_0(x_0+h)h} = -\frac{1}{x_0(x_0+h)} \xrightarrow{h \to 0} -\frac{1}{x_0^2}
\]

(4)
\begin{align*}
\frac{f(x_0+h) - f(x_0)}{h} &= \frac{(x_0+h, x_0^2 + 2x_0h + h^2)^\mathsf{T} - (x_0, x_0^2)^\mathsf{T}}{h} = \frac{(h, 2x_0h + h^2)^\mathsf{T}}{h} \\
&= (1, 2x_0+h)^\mathsf{T} \xrightarrow{h \to 0} (1, 2x_0)^\mathsf{T}
\end{align*}

(5)
\[
\lim_{\substack{n \to \infty \\ n \in \N}} \frac{f(\frac{1}{n}) - f(0)}{\frac{1}{n}} = \frac{\frac{1}{n}}{\frac{1}{n}} = 1; \qquad \lim_{\substack{n \to \infty \\ n \in \N}} \frac{f(-\frac{1}{n}) - f(0)}{-\frac{1}{n}} = \frac{\frac{1}{n}}{-\frac{1}{n}} = -1.
\]
More precisely:
\[
f'(x) = \begin{cases} 1 & \text{for } x > 0 \\ -1 & \text{for } x < 0 \end{cases}
\]

(6) Let $z = x + iy$ with $x, y \in \R$. Then holds:
\[
\frac{f(z+\frac{1}{n}) - f(z)}{\frac{1}{n}} = \frac{x+\frac{1}{n} - x}{\frac{1}{n}} = 1 \quad n \in \N
\]
but:
\[
\frac{f(z+i\frac{1}{n}) - f(z)}{i\frac{1}{n}} = \frac{x - x}{i\frac{1}{n}} = 0.
\]
\hfill $\blacksquare$

\NumberedTheorem{5.1.3}{}{The following statements are equivalent:
\par\medskip
(1) $f$ is differentiable at $x_0$.
\par\medskip
(2) There exists a $c \in Y$ with
\[
\lim_{x \to x_0} \frac{f(x) - f(x_0) - c \cdot (x - x_0)}{x - x_0} = 0.
\]
\par
(3) There exists a $c \in Y$ and a map $r : M \to Y$, which is continuous at $x_0$ and satisfies $r(x_0) = 0$, such that
\[
f(x) = f(x_0) + c(x - x_0) + r(x)(x - x_0) \qquad \forall x \in M
\]
holds.
\par\medskip
In cases (2) and (3), $c$ is the derivative of $f$ at $x_0$.}

\paragraph*{Proof.}
\textbf{(1) $\implies$ (2):} Trivial (Set $c = f'(x_0)$) \\
\textbf{(2) $\implies$ (3):} Define $r$ by
\[
r(x) := \begin{cases} \frac{f(x)-f(x_0)-c(x-x_0)}{x-x_0} & x \neq x_0 \\ 0 & x = x_0 \end{cases}
\]
\[
c := f'(x_0)
\]
\textit{Continuity:} Use (2) and the definition of continuity of its limits. \\
\textbf{(3) $\implies$ (1):} Solve the equation in (3) for $\frac{f(x)-f(x_0)}{x-x_0}$ and use $r(x) \xrightarrow{x \to x_0} 0$. Taking the limit $x \to x_0$ yields $c = f'(x_0)$.
\hfill $\blacksquare$

\NumberedTheorem{5.1.4}{}{If $f$ is differentiable in $x_0$, then $f$ is continuous in $x_0$.}

\paragraph*{Proof.} Use Theorem 5.1.3(3):
\[
f(x) = f(x_0) + \underbrace{c(x - x_0)}_{\xrightarrow{x \to x_0} 0} + \underbrace{r(x)}_{\xrightarrow{x \to x_0} 0} \underbrace{(x - x_0)}_{\xrightarrow{x \to x_0} 0}
\]
From this follows:
\[
\lim_{x \to x_0} f(x) = f(x_0).
\]
\hfill $\blacksquare$

\NumberedRemark{5.1.5}{}{
(1) Ex. 5.1.2(5), (6) show that the converse of Theorem 5.1.4 is generally false.
\par\medskip
(2) A function of the form $x \mapsto A + cx$ is called \textit{affine}. Its graph is a straight line. Theorem 5.1.3(2), (3) states that a function $f$, which is differentiable at $x_0$, can be approximated near $x_0$ by an affine function up to an error $R(x)$ with
\[
\lim_{x \to x_0} \frac{R(x)}{x - x_0} = 0 \qquad (*)
\]
\textbf{Converse:} Let $h : M \to Y$ be an affine function and $R : M \to Y$ a function satisfying $(*)$, as well as
\[
f(x) = h(x) + R(x) \quad \forall x \in M.
\]
Since $h$ is affine, there exists a $c \in Y$ with $h(x) = h(0) + cx$. Then holds:
\[
f(x) - f(x_0) = h(x) + R(x) - h(x_0) - \underbrace{R(x_0)}_{=0} = h(0) + cx + R(x) - h(0) - cx_0 = c(x - x_0) + R(x).
\]
From this follows:
\[
\frac{f(x) - f(x_0)}{x - x_0} = c + \frac{R(x)}{x - x_0} \xrightarrow{x \to x_0} c \qquad \text{because of } (*)
\]
and therefore $f$ is differentiable.
}

\NumberedTheorem{5.1.6}{}{$f = (f_1, \dots, f_m)^T : M \to \K^m$ is differentiable at $x_0 \in M$ if and only if every component function is differentiable at $x_0$. In this case holds:
\[
f'(x_0) = (f_1'(x_0), \dots, f_m'(x_0))^T
\]
}

\paragraph*{Proof.} It holds
\[
\frac{f(x) - f(x_0)}{x - x_0} = \left( \frac{f_1(x) - f_1(x_0)}{x - x_0}, \dots, \frac{f_m(x) - f_m(x_0)}{x - x_0} \right)^T.
\]
A function $f$ is differentiable if the difference quotients $\frac{f(x)-f(x_0)}{x-x_0}$ can be continuously extended for $x \to x_0$. The assertion follows therefore from [8, Theorem 4.3.10] ($g = (g_1, \dots, g_m)^T$ continuous iff all $g_i$ continuous) and [8, Rem. 4.3.13] (continuous extendability). \hfill $\blacksquare$

\NumberedTheorem{5.1.7}{}{Let $f : M \to Y, g : M \to Y$ be differentiable in $x_0$ and $\alpha, \beta \in \K$. Then holds:
\par\medskip
(1) $\alpha f + \beta g$ is differentiable in $x_0$ and it holds
\[
(\alpha f + \beta g)'(x_0) = \alpha f'(x_0) + \beta g'(x_0)
\]
(2) Let $Y = \K$. Then $f \cdot g$ is differentiable in $x_0$ and it holds
\[
(f \cdot g)'(x_0) = f'(x_0)g(x_0) + f(x_0)g'(x_0)
\]
(3) Let $Y = \K$ and $g(x_0) \neq 0$. Then $\frac{f}{g}$ is differentiable in $x_0$ and it holds
\[
\left(\frac{f}{g}\right)'(x_0) = \frac{f'(x_0)g(x_0) - f(x_0)g'(x_0)}{g^2(x_0)}
\]
}

\paragraph*{Proof.}
(1) trivial (Use calculation rules with limits)

(2)
\begin{align*}
\frac{(f \cdot g)(x) - (f \cdot g)(x_0)}{x - x_0} &= \frac{f(x) \cdot g(x) - f(x_0) \cdot g(x_0)}{x - x_0} \\
&= \frac{f(x) - f(x_0)}{x - x_0}g(x) + f(x_0)\frac{g(x) - g(x_0)}{x - x_0} \xrightarrow{x \to x_0} f'(x_0)g(x_0) + f(x_0)g'(x_0).
\end{align*}

(3) Use Theorem 5.1.4. Since therefore $g$ is continuous, there exists a neighborhood $U$ of $x_0$ with $g(x) \neq 0$ for all $x \in U$.
\begin{align*}
\frac{\frac{f(x)}{g(x)} - \frac{f(x_0)}{g(x_0)}}{x - x_0} &= \frac{f(x)g(x_0) - f(x_0)g(x)}{g(x)g(x_0)(x - x_0)} \\
&= \frac{1}{g(x)g(x_0)} \left[ \frac{f(x) - f(x_0)}{x - x_0}g(x_0) - f(x_0)\frac{g(x) - g(x_0)}{x - x_0} \right] \\
&\xrightarrow{x \to x_0, x \in U} \frac{1}{g^2(x_0)} [f'(x_0)g(x_0) - f(x_0)g'(x_0)]
\end{align*}
\hfill $\blacksquare$

\NumberedTheorem{5.1.8}{(Chain Rule)}{Let $f : M \to \K$ be differentiable at $x_0$, $y_0 = f(x_0)$ be an accumulation point of $N = f(M)$ and $g : N \to Y$ be differentiable at $y_0$. Then $g \circ f$ is differentiable at $x_0$ and it holds:
\[
(g \circ f)'(x_0) = g'(f(x_0)) \cdot f'(x_0).
\]
}

\paragraph*{Proof.} Use Theorem 5.1.3 for the characterization of differentiability:

\[
(*) \quad
\begin{aligned}
f(x) &= f(x_0) + f'(x_0)(x - x_0) + r(x)(x - x_0) \\
g(y) &= g(y_0) + g'(y_0)(y - y_0) + s(y)(y - y_0)
\end{aligned}
\qquad r, s \text{ continuous}, r(x_0) = 0, s(y_0) = 0.
\]

\begin{align*}
(g \circ f)(x) &= (g \circ f)(x_0) + g'(f(x_0))(f(x) - f(x_0)) + s(f(x))(f(x) - f(x_0)) \\
&= (g \circ f)(x_0) + g'(f(x_0))f'(x_0)(x - x_0) \\
&\quad + \underbrace{(g'(f(x_0))r(x) + s(f(x))(f'(x_0) + r(x)))}_{=: t(x)}(x - x_0)
\end{align*}

Use [8, Theorem 4.3.5] (Sums, products etc. of continuous functions are continuous) and [8, Theorem 4.3.8] (Composition of two continuous functions is continuous), to show the continuity of $t$. Obviously $t(x_0) = 0$ holds and therefore:
\[
(g \circ f)'(x_0) = g'(f(x_0)) \cdot f(x_0).
\]
\hfill $\blacksquare$

\NumberedTheorem{5.1.9}{(Differentiability of the Inverse Function)}{Let $f : M \to \K$ be injective and differentiable at $x_0$. Let $N = f(M)$. Let $f^{-1} : N \to M$ be continuous at $y_0 := f(x_0)$. Then $f^{-1}$ is differentiable at $y_0$ if and only if $f'(x_0) \neq 0$ holds. In this case:
\[
(f^{-1})'(y_0) = \frac{1}{f'(x_0)}.
\]
}

\paragraph*{Proof.}

(1) Show: $y_0$ is an accumulation point in $N$. $f$ is continuous at $x_0$ (Theorem 5.1.4). Choose $(x_n)$ with $x_n \to x_0, x_n \neq x_0$. Then holds:
\[
f(x_n) \xrightarrow{x_n \to x_0} f(x_0) = y_0
\]
i.e., $y_0$ is an accumulation point.

(2) "$\implies$": For all $x \in M$ holds:
\[
x = f^{-1} \circ f(x)
\]
Differentiating this equation yields:
\[
1 = (f^{-1} \circ f)'(x_0) = (f^{-1})'(f(x_0)) \cdot f'(x_0)
\]
This means: $f'(x_0) \neq 0$ and further:
\[
(f^{-1})'(y_0) = \frac{1}{f'(x_0)}.
\]

"$\Longleftarrow$":
\[
\frac{f^{-1}(y) - f^{-1}(y_0)}{y - y_0} \stackrel{y=f(x), y_0=f(x_0)}{=} \frac{x - x_0}{f(x) - f(x_0)} = \left( \frac{f(x) - f(x_0)}{x - x_0} \right)^{-1} \xrightarrow{x \to x_0} \frac{1}{f'(x_0)}.
\]
Here it was used: $y \neq y_0 \iff x \neq x_0$ (Injectivity of $f$).
\[
y \to y_0 \implies f^{-1}(y) \to f^{-1}(y_0), \quad \text{i.e. } x \to x_0.
\]
\hfill $\blacksquare$

\NumberedRemark{5.1.10}{}{Let $M \subset \R$ be an interval, $M \neq \{x_0\}$. Let $f : M \to \R$ be strictly monotonic, continuous and $f$ differentiable at $x_0$ and $f'(x_0) \neq 0$. Then there exists $f^{-1}$ and is differentiable at $y_0 = f(x_0)$. It holds:
\[
(f^{-1})'(y_0) = \frac{1}{f'(x_0)}
\]
}

\paragraph*{Proof.} Combine Theorem [8, Theorem 4.4.8] with Theorem 5.1.9. \hfill $\blacksquare$

\NumberedExample{5.1.11}{}{$\ln : \R_{>0} \to \R$ is differentiable for all $y_0 \in \R_{>0}$ and it holds:
\[
\ln'(y_0) = \frac{1}{y_0}
\]
}

\paragraph*{Proof.} Set $f : \R \to \R_{>0}, f(x) = e^x$. Then holds:
\[
\ln = f^{-1}.
\]
$f$ satisfies the assumption of Rem. 5.1.10. Then holds:
\[
\ln'(y_0) = (f^{-1})'(y_0) = \frac{1}{f'(x_0)} = \frac{1}{e^{x_0}} = \frac{1}{e^{\ln(y_0)}} = \frac{1}{y_0}, \quad
\left\{
\begin{aligned}
y_0 &= f(x_0), \\
x_0 &= f^{-1}(y_0) = \ln(y_0).
\end{aligned}
\right.
\]
\hfill $\blacksquare$

\NumberedDefinition{5.1.12}{}{A subset $M$ of a normed vector space is called \textit{perfect}, if every point in $M$ is an accumulation point.}

\NumberedExample{5.1.13}{}{
(1) $M$ open $\implies M$ perfect.
\par
(2) $M \subset \R$ interval, $M \neq \emptyset, M \neq \{x_0\}$ is perfect.
\par
(3) $M = \{x_0\}$ is not perfect.
}

\NumberedDefinition{5.1.14}{}{
(1) $f : M \to Y$ is called \textit{differentiable (in $M$)}, if $f$ is differentiable at every point.
\par
(2) $f : M \to Y$ is called \textit{continuously differentiable (in $M$)}, if $f$ is differentiable at every point and $f'$ is continuous in $M$.
\par
(3) $f^{(0)} = f, \, f^{(1)} := Df = f'$, recursively $f^{(n+1)} = (f^{(n)})'$, if $f^{(n)}$ exists and is differentiable.
\par
(4) $f : M \to Y$ is called \textit{$n$-times continuously differentiable}, if $f^{(n)}$ exists and is continuous.
\par
(5) $C^n(M, Y) := \{ f \in C(M, Y) : f \text{ is } n\text{-times continuously differentiable} \}$,
\[
C^\infty(M, Y) := \bigcap_{n \in \N} C^n(M, Y).
\]
}

We now turn to the question, under which conditions a continuous function possesses a maximum and a minimum on its domain.

\NumberedExample{5.1.15}{}{Let $I = ]-1, 1[$ and $f : I \to \R$ be given by $f(x) = x$. Then there exists \textbf{no} $\xi \in I$ and \textbf{no} $\zeta \in I$ with
\[
\forall x \in I : f(x) \le f(\xi),
\]
\[
\forall x \in I : f(x) \ge f(\zeta).
\]
In this example the function $f$ was continuous but the interval open. As we will see in the following, every continuous function on closed intervals assumes its maximum and minimum.
\par
In the following let $(X, \|\cdot\|_X), (Y, \|\cdot\|_Y)$ be normed vector spaces. The set consisting of all subsets of $X$ is called Power set of $X$ and denoted by $\mathcal{P}(X)$.}

\NumberedDefinition{5.1.16}{}{Let $M \subset X$. A set $\{O_\alpha, \alpha \in A\} \subset \mathcal{P}(X)$, consisting of subsets $O_\alpha$ of $X$, is called (open) \textit{Cover} of $M$, if holds:
\[
M \subset \bigcup_{\alpha \in A} O_\alpha \qquad (\text{and } O_\alpha \text{ is open for all } \alpha \in A).
\]
$M$ is called \textit{compact}, if every open cover of $M$ contains a finite subcover, i.e., if there exist an $n \in \N$ and $\alpha_1, \alpha_2, \dots, \alpha_n \in A$ with
\[
M \subset \bigcup_{i=1}^n O_{\alpha_i}.
\]
}

\NumberedExample{5.1.17}{}{
(1) $\emptyset$ is compact.
\par
(2) $\{0\} \cup \{\frac{1}{n} : n \in \N\} \subset \R$ is compact.
\par
(3) $\{\frac{1}{n} : n \in \N\} \subset \R$ is not compact.
}

\NumberedTheorem{5.1.18}{}{Let $M \subset X$ be compact. Then holds:
\par\medskip
(1) $M$ is bounded, i.e. there exists an $R > 0$ with $M \subset B(0, R)$.
\par\medskip
(2) $M$ is closed.
}

\paragraph*{Proof.}
(1) $\{B(0, r) : r \in \R_{\ge 0}\}$ is an open cover of $M$. Using Definition 5.1.16 shows the existence of a finite subcover and therefore an $R > 0$ with $M \subset B(0, R)$.
\par
(2) Show $M^c$ is open. Choose $x \in M^c$ and show: There ex. $r > 0$ with $B(x, r) \subset M^c$.
\par
The system
\[
\left\{ B\left(y, \frac{1}{2}\|y - x\|_X\right) : y \in M \right\}
\]
is an open cover of $M$. Let
\[
\left\{ B\left(y_i, \frac{1}{2}\|y_i - x\|_X\right) : 1 \le i \le n \right\}
\]
be a finite subcover. Define
\[
r := \min_{1 \le i \le n} \frac{1}{2} \|y_i - x\|_X > 0.
\]
Then holds:
\[
B(x, r) \cap B(y_i, r) = \emptyset \quad \forall 1 \le i \le n \implies B(x, r) \cap \bigcup_{i=1}^n B(y_i, r) = \emptyset \implies B(x, r) \cap M = \emptyset.
\]
\hfill $\blacksquare$

\NumberedTheorem{5.1.19}{}{The following statements are equivalent:
\par\medskip
(1) $M \subset X$ is compact.
\par\medskip
(2) Every sequence in $M$ possesses an accumulation point (HP) in $M$.
}

\paragraph*{Proof. $1 \implies 2$:}
Assume (2) were false. Then there exists a sequence $(x_n)_{n \in \N} \subset M$, which possesses no HP in $M$. That is, for every $x \in M$, there is an $\eps(x) > 0$, such that $B(x, \eps(x))$ contains only finitely many sequence terms. $\{B(x, \eps(x)) : x \in M\}$ is an open cover of $M$. Let $\{B(y_i, \eps(y_i)) : 1 \le i \le n\}$ be a finite subcover. Since every $B(y_i, \eps(y_i))$ contains only finitely many sequence terms,
\[
\bigcup_{1 \le i \le n} B(y_i, \eps(y_i))
\]
also contains only finitely many sequence terms and this is a contradiction.

\paragraph*{$2 \implies 1$:}
We show first, that for every $r \in \R_{>0}$ the open cover
\[
B_r := \{B(x, r) : x \in M\}
\]
of $M$ contains a finite subcover.
Assume this is not the case. Then there is an $R > 0$, so that every finite subset of $B_R$ does not cover the set $M$. Let $x_1 \in M$ be arbitrary. Then there is an $x_2 \in M \setminus B(x_1, R)$. Further there is an $x_3 \in M \setminus [B(x_1, R) \cup B(x_2, R)]$. By induction one shows, that there is a sequence $(x_n)_{n \in \N} \subset M$ with
\[
x_{n+1} \in M \setminus \bigcup_{1 \le i \le n} B(x_i, R).
\]
According to (2), $(x_n)_{n \in \N}$ possesses an HP $x^* \in M$. Then $B(x^*, \frac{R}{2})$ contains infinitely many sequence terms $(x_{n_i})_{i \in \N}$. For any two such sequence terms $x_{n_1}$ and $x_{n_2}$ holds
\[
\|x_{n_1} - x_{n_2}\|_X \le \|x_{n_1} - x^*\|_X + \|x_{n_2} - x^*\|_X < R
\]
in contradiction to the construction.

Now let $\{O_\alpha : \alpha \in A\}$ be an arbitrary open cover of $M$. We assume, that it contains no finite subcover. As we have shown above, $B_{\frac{1}{n}}$ contains for every $n \in \N$ a finite subcover. Since $\{O_\alpha : \alpha \in A\}$ contains no finite subcover, there is thus for every $n \in \N$ an $x_n \in M$, so that $B(x_n, \frac{1}{n}) \cap M$ cannot be covered by finitely many $O_\alpha$'s. The sequence $(x_n)_{n \in \N}$ possesses according to (2) an HP $x^*$ in $M$. Then there is an $\alpha^* \in A$ and an $\eps^* \in \R_{>0}$ with
\[
x^* \in O_{\alpha^*} \text{ and } B(x^*, \eps^*) \subset O_{\alpha^*}.
\]
Since $x^*$ is an HP of $(x_n)_{n \in \N}$, there is $N > (\frac{\eps^*}{2})^{-1}$ a $m \ge N$ with
\[
\|x_m - x^*\|_X < \frac{\eps^*}{2}.
\]
For $x \in B(x_m, \frac{1}{m})$ follows
\[
\|x - x^*\|_X \le \|x - x_m\|_X + \|x_m - x^*\| < \frac{1}{m} + \frac{\eps^*}{2} \le \frac{1}{N} + \frac{\eps^*}{2} < \eps^*.
\]
Thus is
\[
B\left(x_m, \frac{1}{m}\right) \subset B(x^*, \eps^*) \subset O_{\alpha^*}
\]
and this is a contradiction. \hfill $\blacksquare$

\NumberedRemark{5.1.20}{}{A set $M$ with property (2) is called \textit{sequentially compact}. Theorem 5.1.19: $X$ norm. VR $\implies$ (``compact'' $\iff$ ``sequentially compact''). But there exist topological VR, for which ``$\iff$'' does not hold.}

\NumberedTheorem{5.1.21}{(Heine-Borel)}{A subset $M \subset \K^m$ is compact if and only if it is bounded and closed.}

\paragraph*{Proof.}
``$\implies$'': Theorem 5.1.18.
\par
``$\Longleftarrow$'':
Let $M \subset \K^m$ be bounded and closed. Let $(x_n)$ be a sequence in $M$. According to Theorem 3.1.15 and Theorem 3.2.5 there exists an $x^* \in \K^m$ and a subsequence $(x_{n_k})_{k \in \N}$ of $(x_n)$ with
\[
(x_{n_k})_i \longrightarrow x^*_i \qquad \text{for } k \to \infty, 1 \le i \le m.
\]
Since $M$ is closed, $x^* \in M$, see Theorem 4.2.14. Therefore $M$ is sequentially compact and because of Theorem 5.1.19 compact.
\hfill $\blacksquare$

\NumberedRemark{5.1.22}{}{The theorem of Heine-Borel does not hold for infinite-dimensional vector spaces. (Ex. $\ell_\infty, \ell_1, \ell_2$)}

\NumberedTheorem{5.1.23}{}{Let $M \subset X, M \neq \emptyset$, compact and $f \in C(M, \R)$. Then $f$ assumes its Maximum and Minimum on $M$, i.e. there exist $\underline{x} \in M$ and $\bar{x} \in M$ with:
\[
f(\underline{x}) \le f(x) \le f(\bar{x}) \qquad \forall x \in M
\]
}

\paragraph*{Proof.} We show first: $M$ compact, $f$ continuous $\implies f(M)$ compact in $\R$. Let $\{O_\alpha : \alpha \in A\}$ be an open cover of $f(M)$. Since $f$ is continuous, we know that $f^{-1}(O_\alpha)$ is relatively open and therefore $\{f^{-1}(O_\alpha) : \alpha \in A\}$ is a relatively open cover of $M$. This possesses a finite subcover $\{f^{-1}(O_{\alpha_i}) : 1 \le i \le n\}$. Therefore $\{O_{\alpha_i} : 1 \le i \le n\}$ is a finite subcover of $f(M)$ and therefore $f(M)$ compact.
\par
Because of Theorem 5.1.18 $f(M)$ is bounded and closed. Thus holds:
\[
-\infty < \rho := \inf f(M) \le \sup f(M) =: R < \infty, \qquad \rho, R \in \R.
\]
From the Def. of the Supremum follows, that there is a sequence $(x_n)$ in $M$ with
\[
\lim_{n \to \infty} f(x_n) = R.
\]
Since $M$ compact, $(x_n)$ possesses an HP $\bar{x}$ in $M$ (Theorem 5.1.19). Since $f$ is continuous, holds
\[
f(\bar{x}) = f(\lim_{k \to \infty} x_{n_k}) = \lim_{k \to \infty} f(x_{n_k}) = R.
\]
The proof of the existence of $\underline{x} \in M$ is analogous. \hfill $\blacksquare$

\NumberedRemark{5.1.24}{}{
(1) $C^n(M, Y)$ is a $\K$-vector space.
\par
(2) $C(M, Y) \supsetneq C^1(M, Y) \supsetneq C^2(M, Y) \supsetneq \dots$
\par
(3) If $M$ is perfect and compact, a norm on $C^n(M, Y)$ is defined by
\[
\|f\|_{C^n(M, Y)} := \max_{0 \le k \le n} \|f^{(k)}\|_{C(M, Y)} = \max_{0 \le k \le n} \max_{x \in M} \|f^{(k)}(x)\|
\]
}

\paragraph*{Proof.}

(1) Follows by induction (analogous to Theorem 4.3.5 ($C(M, Y)$ is a $\K$-VS)) with the help of Theorem 5.1.7(1) ($(\alpha f + \beta g)' = \alpha f' + \beta g'$).
\par
(2) $C^n(M, Y) \subset C^{n-1}(M, Y)$ for all $n \in \N$ follows by induction from Theorem 5.1.3 ($C^1(M, Y) \subset C(M, Y)$).
\par
For $n \ge 1$, $|x|x^{n-1}$ is in $C^{n-1}(\R, \R)$ but not in $C^n(\R, \R)$.
\par
(3) By induction in Exercise.
\hfill $\blacksquare$

\NumberedTheorem{5.1.25}{(Leibniz Rule)}{Let $n \in \N$ and $f, g \in C^n(M, \K)$ with $M \subset \K$ perfect. Then $f \cdot g \in C^n(M, \K)$ and it holds:
\[
(f \cdot g)^{(n)} = \sum_{k=0}^n \binom{n}{k} f^{(k)} \cdot g^{(n-k)}.
\]
}

\paragraph*{Proof.} By induction on $n$. \\
$\mathbf{n = 1}:$ Theorem 5.1.7(2) Product Rule. \\
$\mathbf{n \to n+1}:$
\begin{align*}
(f \cdot g)^{(n+1)} &= ((f \cdot g)^{(n)})' = \sum_{k=0}^n \binom{n}{k} (f^{(k)} \cdot g^{(n-k)})' \\
&= \sum_{k=0}^n \binom{n}{k} (f^{(k+1)} \cdot g^{(n-k)} + f^{(k)} \cdot g^{(n+1-k)}) \\
&= \sum_{l=1}^{n+1} \binom{n}{l-1} f^{(l)} g^{(n+1-l)} + \sum_{k=0}^n \binom{n}{k} f^{(k)} g^{(n+1-k)} \\
&= f^{(n+1)}g + fg^{(n+1)} + \sum_{k=1}^n \underbrace{\left( \binom{n}{k-1} + \binom{n}{k} \right)}_{= \binom{n+1}{k}} f^{(k)} g^{(n+1-k)} \\
&= \sum_{k=0}^{n+1} \binom{n+1}{k} f^{(k)} g^{(n+1-k)}
\end{align*}
\hfill $\blacksquare$

\NumberedExample{5.1.26}{}{
(1) $M \subset \R$, $f : M \to \C$. Then $f = g + ih$, $g : M \to \R$, $h : M \to \R$.
It holds: $f$ differentiable in $x_0 \in M \iff g, h$ differentiable in $x_0$.
\par\medskip
(2) $p(x) = \sum_{k=0}^n a_k x^k, \, a_k \in \K$. Then $p \in C^\infty(\K, \K)$ with
\[
p'(x) = \sum_{k=1}^n k a_k x^{k-1} = \sum_{l=0}^{n-1} (l+1) a_{l+1} x^l
\]
and $p^{(m)}(x) = 0 \quad \forall x \in \K$ for $m > n$.
\par\medskip
(3) $p(x) = \sum_{k=0}^n a_k x^k, \, q(x) = \sum_{k=0}^m b_k x^k$ with $a_k, b_k \in \K$ and $q(x) \neq 0$ for all $x \in M$, $M$ perfect. Then is $\frac{p}{q} \in C^\infty(M, \K)$.
\par\medskip
(4) $\exp \in C^\infty(\K, \K), \, \exp' = \exp$.
\par\medskip
(5) $\sin, \cos \in C^\infty(\K, \K), \, \sin' = \cos, \, \cos' = -\sin$.
\par\medskip
(6) $\ln \in C^\infty(\R_{>0}, \R), \, \ln^{(k)}(x) = \frac{(-1)^{k+1}(k-1)!}{x^k}$ for $k \ge 1$.
\par\medskip
(7) $a > 0$. Then is $x \mapsto a^x \in C^\infty(\R, \R)$ and $(a^x)' = a^x \cdot \ln(a)$.
\par\medskip
(8) $f(x) = \begin{cases} x^2 \cdot \sin \frac{1}{x} & x \in \R_{\neq 0} \\ 0 & x = 0 \end{cases}$ is differentiable on $\R$. The derivative is not continuous in 0.
}

\paragraph*{Proof.}

(1) Follows from differentiability properties of multi-component functions and identification of $\C$ with $\R^2$.
\par
(2) Follows from derivative rules for sums and powers (Theorem 5.1.7, Ex. 5.1.2(1)).
\par
(3) By induction via quotient rule.
\par
(4) Ex. 5.1.2(2).
\par
(5) Use exponential representation of $\sin, \cos$ (Theorem 4.5.2(7)) and (4).
\par
(6) Use $\ln'(x) = \frac{1}{x}$ and complete induction.
\par
(7) Use $a^x = e^{x \cdot \ln a}$, (4), Chain rule.
\par
(8) From $\left| \frac{f(x)-f(0)}{x-0} \right| = \left| \frac{f(x)}{x} \right| = |x| \cdot \underbrace{|\sin \frac{1}{x}|}_{\le 1} \le |x| \xrightarrow{x \to 0} 0$ follows the differentiability in $x = 0$.
\par
For $x \neq 0$ follows from Theorem 5.1.7, 5.1.8 (Chain, Quotient rule) and Part (5)
\[
f'(x) = 2x \sin \frac{1}{x} + x^2 \left( \cos\left(\frac{1}{x}\right) \right) \left(-\frac{1}{x^2}\right) = 2x \sin \frac{1}{x} - \cos \frac{1}{x}.
\]
Thus follows:
\[
f'\left(\frac{1}{2\pi n}\right) = \underbrace{\frac{1}{\pi n} \sin(2\pi n)}_{=0} - \cos(2\pi n) = -1.
\]
\hfill $\blacksquare$

\NumberedDefinition{5.1.27}{}{Let $M \subset \R$ and $x_0$ be an accumulation point of $M$. Then $f$ is called \textit{left-sided} (resp. \textit{right-sided}) \textit{differentiable} at $x_0$, if the limit
\[
(D_-f)(x_0) = \lim_{\substack{x \to x_0 - 0 \\ x \neq x_0}} \frac{f(x) - f(x_0)}{x - x_0} \qquad \left( \text{resp. } (D_+f)(x_0) = \lim_{x \to x_0 + 0} \frac{f(x) - f(x_0)}{x - x_0} \right)
\]
exists.}

\NumberedExample{5.1.28}{}{Let $f(x) = |x|$. Then holds $(D_+f)(0) = 1$ and $(D_-f)(0) = -1$.}

\NumberedTheorem{5.1.29}{}{Let $M \subset \R$. Then $f : M \to Y$ is differentiable at $x_0$ if and only if the left- and right-sided derivatives exist and coincide:
\[
(D_+f)(x_0) = (D_-f)(x_0)
\]
}

\paragraph*{Proof.}
``$\implies$'' trivial. \\
``$\Longleftarrow$'' Define $\phi : M \to Y$ by
\[
\phi(x) := \begin{cases} \frac{f(x)-f(x_0)}{x-x_0} & x \neq x_0 \\ (D_+f)(x_0) & x = x_0 \end{cases}.
\]
$f$ is differentiable at $x_0$ exactly when $\phi$ is continuous at $x_0$. Since $\phi$ is one-sided continuous and $(D_-f)(x_0) = (D_+f)(x_0)$ holds, the assertion follows from Theorem 4.3.17. \hfill $\blacksquare$

\NumberedExample{5.1.30}{}{The function
\[
f(x) := \begin{cases} 0 & x \le 0 \\ e^{-\frac{1}{x}} & x > 0 \end{cases}
\]
is infinitely differentiable and it holds:
\[
f^{(n)}(0) = 0 \qquad \forall n \in \N.
\]
}

\paragraph*{Proof.} Obviously $f \in C^\infty(]-\infty, 0[, \R)$ and $f^{(n)}(x) = 0 \quad \forall n \in \N, x \le 0$. \\
\textbf{Intermediate claim:} $\forall n \in \N \, \exists p_{2n}$ Polynomial of degree $\le 2n$:
\[
f^{(n)}(x) = p_{2n}\left(\frac{1}{x}\right) e^{-\frac{1}{x}} \quad \forall x > 0
\]
\textbf{Proof of the intermediate claim:} $n=0: p_0(y) := 1$. \quad $n \to n+1$:
\[
\left( p_{2n}\left(\frac{1}{x}\right) e^{-\frac{1}{x}} \right)' = \left[ \left( \left(\frac{1}{x^2}\right) p'_{2n}\left(\frac{1}{x}\right) \right) + p_{2n}\left(\frac{1}{x}\right) \left(\frac{1}{x^2}\right) \right] e^{-\frac{1}{x}} =: p_{2n+2}\left(\frac{1}{x}\right) e^{-\frac{1}{x}}.
\]
From this follows $f \in C^\infty(\R_{>0}, \R)$. Because of Theorem 5.1.29 it remains to be shown, that for $n \in \N$ holds:
\[
\lim_{x \to 0+0} f^{(n)}(x) = 0.
\]
From Theorem 4.5.4 ($\frac{y^k}{e^y} \xrightarrow{y \to +\infty} 0$) follows:
\begin{align*}
\lim_{x \to 0+0} f^{(n)}(x) &= \lim_{x \to 0+0} \left( p_{2n}\left(\frac{1}{x}\right) e^{-\frac{1}{x}} \right) = \lim_{x \to 0+0} \sum_{k=0}^{2n} a_k x^{-k} e^{-\frac{1}{x}} \\
&= \sum_{k=0}^{2n} a_k \lim_{x \to 0+0} (x^{-k} e^{-\frac{1}{x}}) = \sum_{k=0}^{2n} a_k \underbrace{\lim_{y \to +\infty} (y^k e^{-y})}_{\to 0} = 0.
\end{align*}
\hfill $\blacksquare$

\end{document}